\PaperTitle{Seeing the Off-Gas Before It Arrives:}{Low-Rank Flame Spectroscopy and Plug-Flow Soft Sensing in Basic Oxygen Furnaces}

In basic oxygen steelmaking, mouth flame radiation responds to bath reactions immediately, whereas flue-gas analyzers report them roughly half a minute later after duct transit. Bridging this temporal mismatch requires a \textbf{Delay-Aware Dual-Modal Framework} that cleans both streams, quantifies variable transport lag, and infers off-gas composition from the flame ahead of the analyzer.

Four pre-modelling checks resolve structural anomalies across the ten heats. Adopting the corrected Attachment 1 restores missing gas records, while oxygen increments and blow duration establish a 1 s sampling step for Heat 5. Because the cumulative-oxygen counter advances in steps of about 4.73 m\ensuremath{^{3}} with a double step every 11 s, a 24-s regression slope replaces direct differencing to recover lance flow. Moreover, carbon–oxygen stoichiometry reveals that reading CO and CO\ensuremath{_{2}} as flow rates would demand \textbf{5.5–6.5 times} the supplied oxygen; we therefore treat the analyzer outputs as volume fractions.

\textbf{For Task 1}, to remove spectral outliers and construct continuous surfaces, a \textbf{Rank-Constrained Robust Low-Rank (RLR)} model decomposes each log spectrum into three physical components and sparse residuals. The leading component tracks overall flame brightness, whereas the second captures spectral colour. In synthetic injection tests, RLR achieves a spike F1 score of \textbf{0.89} against 0.22 for a Hampel filter, while holding out valid pixels yields a \textbf{2.0\%} relative reconstruction error. Sparse outliers localize on the potassium and sodium lines, which vary independently of the thermal continuum.

\textbf{For Task 2}, we reconstruct six continuous process curves per heat and identify four chemical blowing stages with breakpoints at \textbf{23\%, 43\%, and 90\%} of total oxygen (piecewise R\ensuremath{^{2}} = 0.956–0.997). In the cumulative-oxygen coordinate, all ten heats collapse onto one trajectory fitted by the \textbf{Plug-Flow Delay (PFD)} model (R\ensuremath{^{2}} = 0.976). The resulting law, $\tau(Q) = 5 + \frac{127.1}{Q}\,\mathrm{s}$, shortens transport lag from 47.4 s at 3 m\ensuremath{^{3}}/s to \textbf{30.4 s} at 5 m\ensuremath{^{3}}/s. Model-free windowed cross-correlation independently verifies this inverse flow dependence (Spearman \ensuremath{\rho} = \ensuremath{-}0.75, p = 3.4\ensuremath{\times}10\ensuremath{^{-9}}).

\textbf{For Task 3}, a \textbf{Dual-Modal Soft Sensor (DSS)} projects each incoming frame onto a cross-heat low-rank spectral basis, aligns targets via the PFD arrival map, and combines heat-level predictors through causal similarity weighting. Under leave-one-heat-out validation across all ten heats, zero-probe DSS attains mean R\ensuremath{^{2}} values of \textbf{0.873} for CO and \textbf{0.658} for CO\ensuremath{_{2}}. Incorporating Smith-type bias correction with one probe reading every 120 s (\textbf{99.6\% fewer readings}) lifts mean R\ensuremath{^{2}} to \textbf{0.908} and \textbf{0.695}, about 30 s ahead of the analyzer.

Sensitivity checks confirm that perturbing the duct volume by \ensuremath{\pm}10\% shifts zero-probe R\ensuremath{^{2}} by at most 0.003, and that the stage and delay parameters hold across both converter groups. Mouth flame spectra can thus replace nearly all probe readings while delivering off-gas states earlier.

\keywords{robust low-rank decomposition; plug-flow delay; soft sensor; leave-one-heat-out validation; Smith predictor; basic oxygen furnace}


